% Worked example for the Experiment 3 training policies.
%
% Panel A is the real content of training episode train:1:k2, regenerated from
% exp3_training_transfer/coin_city_structural/make_dataset.py (_training_base(1));
% backgrounds are abridged from the rendered prompt.
%
% Panel B draws the prompt-to-key pairing that defines the three trained arms.
% The population-prior control is a cyclic shift within an evidence-depth
% stratum (np.roll under seed 0xC01AC17 in build_training), so it is drawn as a
% shift with one wrap-around edge: every key is still used exactly once, but no
% prompt keeps its own.
\begin{center}
\begin{minipage}{\linewidth}
\captionof{table}{\textbf{One training episode, and what each policy is rewarded against.}
Panel~A is the prompt for episode \texttt{train:1:k2}, abridged from the rendered
text. Panel~B is the prompt--key pairing used during training, drawn for four episodes
at one evidence depth: $P_i$ is episode $i$'s prompt and $\mathbf y^\star_i$ its
ten-answer key. Episode-matched and population-prior training use byte-identical prompts
and the same multiset of keys; only the pairing differs. The control shifts
cyclically across all 1,200 episodes at a depth, so its key is a real key drawn
from the same distribution but never this prompt's, and no reward can be earned by
reading this row's references. Base is untrained and receives no supervision.}
\label{tab:exp3-training-policies}
\scriptsize
\setlength{\tabcolsep}{2.5pt}
\renewcommand{\arraystretch}{1.12}
% The two paper builds have different measures; shrink a panel only if its
% natural width would overrun the text block.
\providecommand{\fitwidth}[1]{%
  \resizebox{\ifdim\width>\linewidth\linewidth\else\width\fi}{!}{#1}}

% ---- Panel A: the shared prompt ------------------------------------------
\noindent\colorbox{black!4}{%
\begin{minipage}{\dimexpr\linewidth-2\fboxsep\relax}%
\fitwidth{%
\begin{tabular}{@{\ }l@{\ \ }l@{\ }}
\multicolumn{2}{@{\ }l@{\ }}{\textbf{A. Shared prompt.} Coin City; resting level
$50.0$, bounded to $[0,100]$; no equation or coefficient is shown.} \\[1.5pt]
\textsc{reference 1} & ``\dots immediate national alerts \dots react
\textbf{strongly}''\quad 14 wk: $-8.0\!\to\!43.8$, $-12.0\!\to\!40.9$, \dots,
$0.0\!\to\!51.6$ \\
\textsc{reference 2} & ``\dots local civic briefings \dots react
\textbf{weakly}''\quad 14 wk: $+8.0\!\to\!52.6$, $+4.0\!\to\!51.3$, \dots,
$-12.0\!\to\!46.9$ \\
\textsc{target} ($k{=}2$) & ``\dots immediate national alerts \dots react
\textbf{strongly}''\quad $-8.0\!\to\!42.3$, $+4.0\!\to\!53.9$ \\
\textsc{query} & ten counterfactuals A--J: shock $\in\{-12,-6,0,6,12\}$ $\times$
horizon $\in\{1,3\}$; return \texttt{\{"forecasts":[A,\dots,J]\}} \\
\end{tabular}}%
\end{minipage}}

\vspace{6pt}
\noindent\textbf{B. How prompts are paired with answer keys during training.}
\vspace{2pt}

% ---- Panel B: the prompt-to-key wiring ------------------------------------
\definecolor{wireInk}{HTML}{3F4550}
\definecolor{wireMute}{HTML}{8A93A3}
\definecolor{wireKeyFill}{HTML}{E8F2FB}
\definecolor{wireKeyEdge}{HTML}{4F7EA8}
\noindent\fitwidth{%
\begin{tikzpicture}[
  x=1cm, y=1cm,
  pchip/.style={rounded corners=2pt, minimum width=0.95cm, minimum height=0.26cm,
    inner sep=0pt, draw=wireMute!60, fill=black!4, line width=.5pt,
    font=\scriptsize, text=wireInk},
  schip/.style={pchip, dash pattern=on 1.3pt off 1.1pt},
  kchip/.style={rounded corners=2pt, minimum width=0.95cm, minimum height=0.26cm,
    inner sep=0pt, draw=wireKeyEdge, fill=wireKeyFill, line width=.5pt,
    font=\scriptsize, text=wireInk},
  flatbox/.style={kchip, minimum width=1.60cm, minimum height=0.30cm,
    draw=wireMute!70, fill=black!4},
  wire/.style={-{Stealth[length=3pt,width=2.4pt]}, draw=wireInk!65,
    line width=.5pt, shorten >=.5pt, shorten <=.5pt},
  ptitle/.style={anchor=south, font=\scriptsize\bfseries, text=wireInk},
  pnote/.style={anchor=north, font=\tiny\itshape, text=wireMute},
]
\node[ptitle] at (1.55,0.12) {Episode-matched};
\node[ptitle] at (6.65,0.12) {Population prior};
\node[ptitle] at (11.72,0.12) {Structureless};

% --- Episode-matched: each prompt keeps its own key --------------------------------
\foreach \r/\y in {1/-0.30, 2/-0.66, 3/-1.02, 4/-1.38}{
  \node[pchip] (cp\r) at (0.50,\y) {$P_\r$};
  \node[kchip] (ck\r) at (2.60,\y) {$\mathbf y^\star_\r$};
  \draw[wire] (cp\r.east) -- (ck\r.west);}
\node[pnote] at (1.55,-1.85) {its own key};

% --- Population prior: cyclic shift within an evidence depth ---------------
% The wrap edge is routed under the chips and up outside the key column so it
% cannot clip a chip it does not connect to.
\foreach \r/\y in {1/-0.30, 2/-0.66, 3/-1.02, 4/-1.38}{
  \node[pchip] (pp\r) at (5.60,\y) {$P_\r$};
  \node[kchip] (pk\r) at (7.70,\y) {$\mathbf y^\star_\r$};}
\foreach \a/\b in {1/2, 2/3, 3/4}{
  \draw[wire] (pp\a.east) -- (pk\b.west);}
\draw[wire, rounded corners=3.5pt]
  (pp4.south) -- (5.60,-1.68) -- (8.50,-1.68) -- (8.50,-0.30) -- (pk1.east);
\node[pnote] at (6.65,-1.85) {every key used once, none its own};

% --- Structureless: scrambled prompt, one flat key -------------------------
\foreach \r/\y in {1/-0.30, 2/-0.66, 3/-1.02, 4/-1.38}{
  \node[schip] (sp\r) at (10.70,\y) {$\tilde P_\r$};}
\node[flatbox] (sflat) at (12.75,-0.84) {$(50,\dots,50)$};
\foreach \r in {1,2,3,4}{\draw[wire] (sp\r.east) -- (sflat.west);}
\node[pnote] at (11.72,-1.85) {flat key; displayed drivers permuted};
\end{tikzpicture}}
\end{minipage}
\end{center}
